\documentclass[10pt,a4paper]{amsart}
\usepackage[margin=19mm]{geometry}
\usepackage{amsmath,amssymb,mathtools}
\usepackage[scr=rsfs,cal=euler]{mathalfa}
\usepackage{xfrac}
\usepackage[unicode,hidelinks,bookmarks=false]{hyperref}
\newcommand{\ie}{i.e.\ }
\newcommand{\cf}{cf.\ }
\newcommand{\resp}{respectively\ }
\newcommand{\Subsection}{\subsection}
\providecommand{\nobreakdash}{\nobreak\hbox{-}}
\setlength{\emergencystretch}{2em}
\multlinegap0pt
\usepackage{bm}
\usepackage{bbm}
\usepackage{dsfont}
\usepackage{xspace}
\usepackage{cancel}
\usepackage{mathtools}
\usepackage{amsthm}
\makeatletter
\@ifundefined{theorem}{\newtheorem{theorem}{Theorem}}{}
\@ifundefined{claim}{\newtheorem{claim}[theorem]{Claim}}{}
\@ifundefined{lemma}{\newtheorem{lemma}[theorem]{Lemma}}{}
\makeatother
\title[Two-cardinal Kurepa Hypotheses]{Two-cardinal Kurepa Hypotheses}
\author{Fanxin Wu}
\date{Source: arXiv:2510.08860v2 (CC BY 4.0)}
\begin{document}
\maketitle

\noindent\emph{Source and licence.} Two-cardinal Kurepa Hypotheses. Version arXiv:2510.08860v2 (CC BY 4.0), distributed under the Creative Commons Attribution 4.0 International licence, as recorded in the arXiv licence metadata for this exact version and shown on the abstract page https://arxiv.org/abs/2510.08860v2. Full rights notice: licenses/arxiv-2510.08860-CC-BY-4.0.txt.\par\smallskip

\noindent\textit{Editorial scope.} Counted unit: the Theorem whose source label is \texttt{None} in the article (source0.tex lines 709--711, source label \texttt{None}), delivered together with the article's own complete original proof of it (source0.tex lines 712--772, one \texttt{proof} environment of 61 lines). No mathematics is written, repaired or re-derived: statement and proof are the article's own text line for line. Required context retained uncounted: forcing\_Easton (lines 181--187); forcing\_branch\_Unger\_formerly\_closed (lines 226--228); KH\_upward (lines 300--301); KH\_Q\_effect (lines 524--526). Every definition, display and label of those lines is retained unchanged. Omitted from this package and not redistributed: 1--180, 188--225, 229--299, 302--523, 527--708, 773--816. Nothing inside a retained excerpt is omitted or altered: every retained body is delivered exactly as the article's lines are written, so no omission receipt of any kind accompanies this package. Citations: the retained excerpts carry no citation command at all, so no bibliography entry of the article is transcribed and no reference is redistributed. Reference adaptation: none was needed and none was made; every reference inside the delivered unit resolves inside it, so no unresolved reference or citation is produced. Printed slips kept literal: no printed word, symbol, hypothesis or number of the counted statement or of its proof was altered. Layout and class: the article's own class is \texttt{amsart} with packages including amsmath, amssymb, mathrsfs, xcolor, enumitem, comment, todonotes, xy, url, amsthm, thmtools, thm-restate, hyperref, soul; the wrapper is the lane's pinned \texttt{amsart} editorial wrapper at 10pt on A4, which restates the theorem-like environments the retained text needs and adds the article's own macro definitions that the retained text needs (none), with extra wrapper packages (bm, bbm, dsfont, xspace, cancel, mathtools). The delivered numbering is the unit's own. Assets: no retained span contains a figure environment, a graphics-inclusion command, a PDF-page inclusion, a table or a TikZ picture; no image file is copied into this package, the package declares no asset dependency and no image file is redistributed with it. Licence: version arXiv:2510.08860v2 of this article is distributed under the Creative Commons Attribution 4.0 International licence, as recorded in the arXiv licence metadata for this exact version and shown on the abstract page https://arxiv.org/abs/2510.08860v2. A full rights notice is delivered at licenses/arxiv-2510.08860-CC-BY-4.0.txt. Characters: every retained excerpt is pure ASCII, so the delivered engine, XeLaTeX with its default Latin Modern text font, reports no missing character.
\begin{lemma}[Easton]\label{forcing_Easton}
Assume $\kappa$ is a regular cardinal, $\mathbb{P}$ and $\mathbb{Q}$ are forcing posets, $\mathbb{P}$ is $\kappa$-cc and $\mathbb{Q}$ is $\kappa$-closed. Suppose $P$ is $\mathbb{P}$-generic over $V$ and $Q$ is $\mathbb{Q}$-generic over $V$.

(i) $\mathbb{P}$ is $\kappa$-cc in $V[Q]$. Therefore, any antichain of $\mathbb{P}$ in $V[Q]$ is already in $V$, and $P$ is generic over $V[Q]$.

(ii) $\mathbb{Q}$ is $\kappa$-distributive in $V[P]$, and $Q$ is generic over $V[P]$.
\end{lemma}

\begin{lemma}\label{forcing_branch_Unger_formerly_closed}
Suppose $\kappa\leq\lambda$ are regular, $\mathbb{P}$ is $\kappa$-cc, $\mathbb{Q}$ is $\kappa$-closed, and $\exists\tau<\kappa,\ 2^\tau\geq\lambda$. If $T$ is a slim $\lambda$-tree in $V[P]$, then $T$ does not gain any new branch in $V[P][Q]$.
\end{lemma}

\begin{lemma}\label{KH_upward}
Suppose $V\subseteq W$ are transitive models of $\mathsf{ZFC}$, $\kappa$ is a regular cardinal in $V$, and $\mathcal{P}_\kappa(\lambda)^V$ is cofinal in $\mathcal{P}_\kappa(\lambda)^W$. Then $V\models\mathsf{KH}(\kappa,\lambda,\mu)$ implies $W\models\mathsf{KH}(\kappa,\lambda,\mu)$. In particular, this is true if $W=V[P]$ where $\mathbb{P}$ is either $\kappa$-distributive or $\kappa$-cc.\end{lemma}

\begin{lemma}\label{KH_Q_effect}
$\mathbb{Q}(\kappa^+,\lambda,\mu)$ forces $\mathsf{KH}(\kappa^+,\lambda,\mu)$. In fact it forces $\mathsf{KH}(\kappa^+,\theta,\mu)$ for any $\kappa^+\leq\theta\leq\lambda$, so in particular $2^{\kappa^+}\geq\mu$.
\end{lemma}

\begin{theorem}
$\mathsf{TP(\kappa^{++})\land\mathsf{KH}(\kappa^+)}$ for singular $\kappa$ is consistent relative to large cardinals.
\end{theorem}

\begin{proof}
Let $\kappa$ be indestructibly supercompact, $\lambda>\kappa$ be measurable (this is for simplicity; weakly compact should suffice), $\mathbb{M}=\mathbb{M}(\kappa,\lambda)$, $\mathbb{P}\times\mathbb{R}\twoheadrightarrow\mathbb{M}$ be the usual projection, and $\mathbb{Q}=\mathbb{Q}(\kappa^+,\kappa^+,\lambda)$. Following Cummings, we force with $\mathbb{M}\times\mathbb{Q}$, and then use Prikry forcing to singularize $\kappa$. 

\begin{claim}
Working in $V[M\times Q]$, if $\overline{\mathbb{P}}$ is any $\kappa^+$-cc forcing in the Cohen extension $V[P]$, then $V[M\times Q][\overline{P}]\models\mathsf{TP(\kappa^{++})\land\mathsf{KH}(\kappa^+)}$.
\end{claim}


This in particular applies to the Prikry forcing $\mathbb{\overline{P}}$ as defined in $V[P]$ with respect to some normal measure $U$, and thus would finish the proof. Note that $U$ remains a normal measure in $V[P\times R\times Q]$ and thus in $V[M\times Q]$, because $\mathbb{R}$ and $\mathbb{Q}$ are $\kappa^+$-closed in $V$, so by Easton's Lemma \ref{forcing_Easton} the extension $V[P]\subseteq V[P\times R\times Q]$ does not add $<\kappa^+$-sequences.

We first deal with $\mathsf{KH}(\kappa^+)$, which is the easier part. By Lemma \ref{KH_Q_effect} we have $\mathsf{KH}(\kappa^+,\kappa^+,\lambda)$ in $V[Q]$. We have the following sequence of extensions:

\begin{center}
$V[Q]\subseteq V[R\times Q]\subseteq V[P\times R\times Q]$
\end{center}

The first extension is $\kappa^+$-closed, and the second extension is $\kappa^+$-cc by Easton's Lemma \ref{forcing_Easton}, so by Lemma \ref{KH_upward} we have $\mathsf{KH}(\kappa^+,\kappa^+,\lambda)$ in $V[P\times R\times Q]$, and hence in $V[M\times Q]$ since it contains the witnessing family. Finally $\overline{\mathbb{P}}$ is $\kappa^+$-cc in $V[M\times Q]$, again by Easton, so $\mathsf{KH}(\kappa^+,\kappa^+,\lambda)$ holds in $V[M\times Q][\overline{P}]$.

Next we argue that there is no $\lambda$-Aronszajn tree. Start with an elementary embedding $j:V\rightarrow W$ with critical point $\lambda$. Let $\mathbb{P}^\lambda\times\mathbb{R}^\lambda\twoheadrightarrow j(\mathbb{M})/M$ be the projection in $W[M]$. There is also a projection $j(\mathbb{Q})\twoheadrightarrow\mathbb{Q}$ in $V$ since $\mathbb{Q}=j(\mathbb{Q})_\lambda$. We lift the embedding $j$ successively as follows.

First we force with $\mathbb{P}^\lambda\times \mathbb{R}^\lambda\times j(\mathbb{Q})/Q$ over $V[M\times Q]$. In the resulting model $V[M][P^\lambda\times R^\lambda][j(Q)]$, we can lift the embedding $j:V\rightarrow W$ to

\begin{center}
$j:V[M\times Q]\rightarrow W[j(M)\times j(Q)]$
\end{center}

\noindent where $j(M)$ is the $j(\mathbb{M})$-generic filter induced by $M*(P^\lambda\times R^\lambda)$; it can be checked that $j(M)$ and $j(Q)$ are mutually generic, so the model $W[j(M)\times j(Q)]$ makes sense.

\iffalse
$V[M][Q]=V[Q][M]\subseteq V[Q][M][j(Q)/Q]\subseteq V[Q][M][j(Q)/Q][P^\lambda\times R^\lambda]= V[Q][j(Q)/Q][M][P^\lambda\times R^\lambda]\supseteq V[j(Q)][j(M)]$
\fi

We can now make sense of $j(\mathbb{\overline{P}})$. Note that $j$ restricts to a map from $\overline{\mathbb{P}}$ into $j(\overline{\mathbb{P}})$ that is almost a complete embedding. More precisely, since $V[M\times Q]\models\mathbb{\overline{P}}$ is $\kappa^+$-cc and $\mathrm{crit}(j)=\lambda>\kappa$, if $A\subseteq\overline{\mathbb{P}}$ is a maximal antichain in $V[M\times Q]$ then $j(A)=j[A]$, so $j[A]$ is also maximal.

It follows that if we let $\overline{P}^*$ be $j(\mathbb{\overline{P}})$-generic over $V[M][P^\lambda\times R^\lambda][j(Q)]$, then its preimage under $j$, denoted $\overline{P}$, is $\overline{\mathbb{P}}$-generic over $V[M\times Q]$. Thus we can lift the embedding to

\begin{center}
$j:V[M\times Q][\overline{P}]\rightarrow W[j(M)\times j(Q)][j(\overline{P})]$
\end{center}

\noindent in the model $V[M][P^\lambda\times R^\lambda][j(Q)][j(\overline{P})]$, where $j(\overline{P})=\overline{P}^*$. Let $T$ be a $\lambda$-tree in $V[M\times Q][\overline{P}]$; by standard arguments we may assume it is in $W[M\times Q][\overline{P}]$. The elementary embedding provides a branch of $T$ in $W[j(M)\times j(Q)][j(\overline{P})]$, and we want to show that the branch is already in $W[M\times Q][\overline{P}]$. We have the following sequence of extensions:

\begin{align}
W[M][Q][\overline{P}]&\subseteq W[M][j(Q)][\overline{P}]\\
&\subseteq W[M][P^\lambda][j(Q)][\overline{P}]\\
&\subseteq W[M][P^\lambda][j(Q)][j(\overline{P})]\\
&\subseteq W[M][P^\lambda\times R^\lambda][j(Q)][j(\overline{P})]
\end{align}

It suffices to argue that none of the extension adds branch to $T$.

\iffalse$\overline{P}$ is generic over $V[M\times Q]$, so in particular over $W[M]$. By successive use of Easton, $\overline{\mathbb{P}}$ remains $\kappa^+$-cc in $W[M][j(Q)]$ and $j(\mathbb{Q})$ is $\kappa^+$-distributive in $W[M]$, so it does not add antichains in $\overline{\mathbb{P}}$. Thus $\overline{P}$ is generic over $W[M][j(Q)]$, and we are free to interchange the order of $\overline{P}$ and $j(Q)$.\fi

(1) $j(\mathbb{Q})/Q$ is square-$\lambda$-cc in $W[Q]$ and thus in $W[M][Q][\overline{P}]$, because $\mathbb{M}$ and $\overline{\mathbb{P}}$ are both $\lambda$-Knaster. Thus the first extension does not add a new branch to $T$.

(2) $\mathbb{P}^\lambda$ is $\kappa^+$-Knaster in $W[M][j(Q)]$ where it is still Cohen forcing, and thus at least $\kappa^+$-cc in $W[M][j(Q)][\overline{P}]$ because $\mathbb{\overline{P}}$ is $\kappa^+$-cc. Since $\kappa^+<\lambda$, the second extension does not add branch.

(3) By elementarity $j(\mathbb{\overline{P}})/\overline{P}$ is $\kappa^+$-cc in $W[j(M)\times j(Q)][\overline{P}]$, and thus in the smaller model $W[M][P^\lambda][j(Q)][\overline{P}]$, so the third extension does not add branch.

(4) $\mathbb{R}^\lambda$ is $\lambda$-closed in $W[M]$. Applying Lemma \ref{forcing_branch_Unger_formerly_closed} in $W[M]$ to $(\mathbb{P}^\lambda\times j(\mathbb{Q})) *j(\mathbb{\overline{P}})$ and $\mathbb{R}^\lambda$, we see that the fourth extension does not add branch. Note that technically $j(\mathbb{\overline{P}})$ is a $j(\mathbb{P})$-name in $V$ but we are viewing it as a $\mathbb{P}^\lambda$-name in $V[P]$.
\end{proof}

\end{document}
